% GENERATED FILE. Edit canonical lessons, then run npm run generate:books.
% canonical-source-hash: fnv1a64:607e817e96f1d868
% canonical-lessons: LA-C83-culina, LA-C83-tectum, LA-C83-lucerna, LA-C83-poculum, LA-C83-culter

\chapter{Kitchen, Roof, Lamp, Cup, Knife}
\label{ch:la-kitchen-roof-lamp-cup-knife}
\hlchaptermodality{83}

\begin{chapteropening}
\textbf{By the end of this chapter:} \emph{I can name a kitchen, a roof, a lamp, a cup and a knife.}

Complete the last lesson of chapter 83: I can name a kitchen, a roof, a lamp, a cup and a knife.
\end{chapteropening}

\section[culīna, culīnae]{culīna, culīnae — a kitchen}

\label{lesson:LA-C83-culina}



\begin{quote}
Before the new one: say the Latin for a window, then the Latin for a bedroom.
\end{quote}

\subsection*{You'll want to know: culīna, culīnae}
\textbf{culīna, culīnae} — \textquotedblleft{}a kitchen\textquotedblright{}. English \emph{culinary}.

\begin{cousinweb}
\textbf{culīna} has no sure origin; a link with \emph{coquō}, \textquotedblleft{}to cook\textquotedblright{}, is often suggested but not proven.
\end{cousinweb}

\begin{grammarlens}[title={What you've built}]
One more word for this chapter.
\end{grammarlens}

\subsection*{Your turn}
\begin{itemize}
\raggedright
  \item \emph{Say it:} \emph{culīna}
  \item \emph{Say it:} \emph{culīna}, once more
  \item \emph{Say it:} say \emph{culīna}
  \item \emph{Recall:} say the Latin for a window, then the Latin for a bedroom, then say \emph{culīna} again
\end{itemize}

\begin{tcolorbox}[breakable,colback=teal!4,colframe=teal!35!black,title={Before you move on}]
What does culīna, culīnae mean? (\textquotedblleft{}A kitchen\textquotedblright{}.) Say it once more.
\end{tcolorbox}
\section[tēctum, tēctī]{tēctum, tēctī — a roof}

\label{lesson:LA-C83-tectum}



\begin{quote}
Before the new one: say the Latin for a bedroom, then the Latin for a kitchen.
\end{quote}

\subsection*{You'll want to know: tēctum, tēctī}
\textbf{tēctum, tēctī} — \textquotedblleft{}a roof\textquotedblright{}.

\begin{cousinweb}
\textbf{tēctum} is the past participle of \emph{tegō}, \textquotedblleft{}to cover\textquotedblright{}, used as a noun: \textquotedblleft{}what covers\textquotedblright{} (Lewis \& Short). \emph{tegō} is usually traced to Proto-Indo-European \emph{*(s)teg-}, \textquotedblleft{}to cover with a roof\textquotedblright{} (EtymDB 2.1); English \emph{thatch} is a cousin.
\end{cousinweb}

\begin{grammarlens}[title={What you've built}]
One more word for this chapter.
\end{grammarlens}

\subsection*{Your turn}
\begin{itemize}
\raggedright
  \item \emph{Say it:} \emph{tēctum}
  \item \emph{Say it:} \emph{tēctum}, once more
  \item \emph{Say it:} say \emph{tēctum}
  \item \emph{Recall:} say the Latin for a bedroom, then the Latin for a kitchen, then say \emph{tēctum} again
\end{itemize}

\begin{tcolorbox}[breakable,colback=teal!4,colframe=teal!35!black,title={Before you move on}]
What does tēctum, tēctī mean? (\textquotedblleft{}A roof\textquotedblright{}.) Say it once more.
\end{tcolorbox}
\section[lucerna, lucernae]{lucerna, lucernae — a lamp}

\label{lesson:LA-C83-lucerna}



\begin{quote}
Before the new one: say the Latin for a kitchen, then the Latin for a roof.
\end{quote}

\subsection*{You'll want to know: lucerna, lucernae}
\textbf{lucerna, lucernae} — \textquotedblleft{}a lamp\textquotedblright{}. It gives \textbf{lūx}.

\begin{cousinweb}
\textbf{lucerna} is built on \emph{lūceō}, \textquotedblleft{}to shine\textquotedblright{} (Lewis \& Short), with \emph{-erna}, a suffix also seen in \emph{caverna}, \textquotedblleft{}a hollow\textquotedblright{}. \emph{lūceō} belongs with \emph{lūx}, \textquotedblleft{}light\textquotedblright{}.
\end{cousinweb}

\begin{grammarlens}[title={What you've built}]
One more word for this chapter.
\end{grammarlens}

\subsection*{Your turn}
\begin{itemize}
\raggedright
  \item \emph{Say it:} \emph{lucerna}
  \item \emph{Say it:} \emph{lucerna}, once more
  \item \emph{Say it:} say \emph{lucerna}
  \item \emph{Recall:} say the Latin for a kitchen, then the Latin for a roof, then say \emph{lucerna} again
\end{itemize}

\begin{tcolorbox}[breakable,colback=teal!4,colframe=teal!35!black,title={Before you move on}]
What does lucerna, lucernae mean? (\textquotedblleft{}A lamp\textquotedblright{}.) Say it once more.
\end{tcolorbox}
\section[pōculum, pōculī]{pōculum, pōculī — a cup}

\label{lesson:LA-C83-poculum}



\begin{quote}
Before the new one: say the Latin for a roof, then the Latin for a lamp.
\end{quote}

\subsection*{You'll want to know: pōculum, pōculī}
\textbf{pōculum, pōculī} — \textquotedblleft{}a cup\textquotedblright{}.

\begin{cousinweb}
\textbf{pōculum} is inherited, usually traced through Proto-Italic \emph{*pōtlom}, \textquotedblleft{}drinking cup\textquotedblright{}, to Proto-Indo-European \emph{*peh3tlom}, \textquotedblleft{}a thing to drink from\textquotedblright{} (EtymDB 2.1). \emph{pōtō}, \textquotedblleft{}to drink\textquotedblright{}, belongs to the same family.
\end{cousinweb}

\begin{grammarlens}[title={What you've built}]
One more word for this chapter.
\end{grammarlens}

\subsection*{Your turn}
\begin{itemize}
\raggedright
  \item \emph{Say it:} \emph{pōculum}
  \item \emph{Say it:} \emph{pōculum}, once more
  \item \emph{Say it:} say \emph{pōculum}
  \item \emph{Recall:} say the Latin for a roof, then the Latin for a lamp, then say \emph{pōculum} again
\end{itemize}

\begin{tcolorbox}[breakable,colback=teal!4,colframe=teal!35!black,title={Before you move on}]
What does pōculum, pōculī mean? (\textquotedblleft{}A cup\textquotedblright{}.) Say it once more.
\end{tcolorbox}
\section[culter, cultrī]{culter, cultrī — a knife}

\label{lesson:LA-C83-culter}



\begin{quote}
Before the new one: say the Latin for a lamp, then the Latin for a cup.
\end{quote}

\subsection*{You'll want to know: culter, cultrī}
\textbf{culter, cultrī} — \textquotedblleft{}a knife\textquotedblright{}.

\begin{cousinweb}
\textbf{culter} is probably inherited, from a Proto-Indo-European \emph{*kelHtro-}, \textquotedblleft{}knife\textquotedblright{}, built on \emph{*(s)kelH-}, \textquotedblleft{}to split, to cut\textquotedblright{} (EtymDB 2.1). English \emph{coulter}, the blade in front of a ploughshare, comes from \emph{culter}.
\end{cousinweb}

\begin{grammarlens}[title={What you've built}]
One more word for this chapter.
\end{grammarlens}

\subsection*{Your turn}
\begin{itemize}
\raggedright
  \item \emph{Say it:} \emph{culter}
  \item \emph{Say it:} \emph{culter}, once more
  \item \emph{Say it:} say \emph{culter}
  \item \emph{Recall:} say the Latin for a lamp, then the Latin for a cup, then say \emph{culter} again
\end{itemize}

\begin{tcolorbox}[breakable,colback=teal!4,colframe=teal!35!black,title={Before you move on}]
What does culter, cultrī mean? (\textquotedblleft{}A knife\textquotedblright{}.) Say it once more.
\end{tcolorbox}
